\subsection{\label{sec.perm.sign}Signs of permutations}

The notion of the \emph{sign} (aka \emph{signature}) of a permutation is a
simple consequence of that of its length; moreover, it is rather well-known,
due to its role in the definition of a determinant. Thus we will survey its
properties quickly and without proofs. More can be found in \cite[\S 5.3 and
\S 5.6]{detnotes} and \cite[Appendix B]{Strick13}.

\begin{definition}
\label{def.perm.sign}Let $n\in\mathbb{N}$. The \emph{sign} of a permutation
$\sigma\in S_{n}$ is defined to be the integer $\left(  -1\right)
^{\ell\left(  \sigma\right)  }$.

It is denoted by $\left(  -1\right)  ^{\sigma}$ or $\operatorname*{sgn}\left(
\sigma\right)  $ or $\operatorname*{sign}\left(  \sigma\right)  $ or
$\varepsilon\left(  \sigma\right)  $. It is also known as the \emph{signature}
of $\sigma$.
\end{definition}

\begin{proposition}
\label{prop.perm.sign.props}Let $n\in\mathbb{N}$. \medskip

\textbf{(a)} The sign of the permutation $\operatorname*{id}\in S_{n}$ is
$\left(  -1\right)  ^{\operatorname*{id}}=1$. \medskip

\textbf{(b)} For any two distinct elements $i$ and $j$ of $\left[  n\right]
$, the transposition $t_{i,j}\in S_{n}$ has sign $\left(  -1\right)
^{t_{i,j}}=-1$. \medskip

\textbf{(c)} For any positive integer $k$ and any distinct elements
$i_{1},i_{2},\ldots,i_{k}\in\left[  n\right]  $, the $k$-cycle
$\operatorname*{cyc}\nolimits_{i_{1},i_{2},\ldots,i_{k}}$ has sign $\left(
-1\right)  ^{\operatorname*{cyc}\nolimits_{i_{1},i_{2},\ldots,i_{k}}}=\left(
-1\right)  ^{k-1}$. \medskip

\textbf{(d)} We have $\left(  -1\right)  ^{\sigma\tau}=\left(  -1\right)
^{\sigma}\cdot\left(  -1\right)  ^{\tau}$ for any $\sigma\in S_{n}$ and
$\tau\in S_{n}$. \medskip

\textbf{(e)} We have $\left(  -1\right)  ^{\sigma_{1}\sigma_{2}\cdots
\sigma_{p}}=\left(  -1\right)  ^{\sigma_{1}}\left(  -1\right)  ^{\sigma_{2}%
}\cdots\left(  -1\right)  ^{\sigma_{p}}$ for any $\sigma_{1},\sigma_{2}%
,\ldots,\sigma_{p}\in S_{n}$. \medskip

\textbf{(f)} We have $\left(  -1\right)  ^{\sigma^{-1}}=\left(  -1\right)
^{\sigma}$ for any $\sigma\in S_{n}$. (The left hand side here has to be
understood as $\left(  -1\right)  ^{\left(  \sigma^{-1}\right)  }$.) \medskip

\textbf{(g)} We have%
\[
\left(  -1\right)  ^{\sigma}=\prod_{1\leq i<j\leq n}\dfrac{\sigma\left(
i\right)  -\sigma\left(  j\right)  }{i-j}\ \ \ \ \ \ \ \ \ \ \text{for each
}\sigma\in S_{n}.
\]
(The product sign \textquotedblleft$\prod_{1\leq i<j\leq n}$\textquotedblright%
\ means a product over all pairs $\left(  i,j\right)  $ of integers satisfying
$1\leq i<j\leq n$. There are $\dbinom{n}{2}$ such pairs.) \medskip

\textbf{(h)} If $x_{1},x_{2},\ldots,x_{n}$ are any elements of some
commutative ring, and if $\sigma\in S_{n}$, then%
\[
\prod_{1\leq i<j\leq n}\left(  x_{\sigma\left(  i\right)  }-x_{\sigma\left(
j\right)  }\right)  =\left(  -1\right)  ^{\sigma}\prod_{1\leq i<j\leq
n}\left(  x_{i}-x_{j}\right)  .
\]

\end{proposition}

\begin{proof}
[Proof of Proposition \ref{prop.perm.sign.props} (sketched).]Most of this
follows easily from what we have proved above, but here are references to
complete proofs: \medskip

\textbf{(a)} This is \cite[Proposition 5.15 \textbf{(a)}]{detnotes}, and
follows easily from $\ell\left(  \operatorname*{id}\right)  =0$. \medskip

\textbf{(d)} This is \cite[Proposition 5.15 \textbf{(c)}]{detnotes}, and
follows easily from Corollary \ref{cor.perm.red.sigtau} \textbf{(a)}. A
different proof appears in \cite[Proposition B.13]{Strick13}. \medskip

\textbf{(b)} This is \cite[Exercise 5.10 \textbf{(b)}]{detnotes}, and follows
easily from Exercise \ref{exe.perm.len.cycles} \textbf{(a)}. \medskip

\textbf{(c)} This is \cite[Exercise 5.17 \textbf{(d)}]{detnotes}, and follows
easily from Exercise \ref{exe.perm.cycles.c-t} \textbf{(a)} and Exercise
\ref{exe.perm.len.cycles} \textbf{(b)} using Proposition
\ref{prop.perm.sign.props} \textbf{(d)}. \medskip

\textbf{(e)} This is \cite[Proposition 5.28]{detnotes}, and follows by
induction from Proposition \ref{prop.perm.sign.props} \textbf{(d)}. \medskip

\textbf{(f)} This is \cite[Proposition 5.15 \textbf{(d)}]{detnotes}, and
follows easily from Proposition \ref{prop.perm.sign.props} \textbf{(d)} or
from Proposition \ref{prop.perm.len.inv}. \medskip

\textbf{(h)} This is \cite[Exercise 5.13 \textbf{(a)}]{detnotes} (or, rather,
the straightforward generalization of \cite[Exercise 5.13 \textbf{(a)}%
]{detnotes} to arbitrary commutative rings). The proof is fairly easy: Each
factor $x_{\sigma\left(  i\right)  }-x_{\sigma\left(  j\right)  }$ on the left
hand side appears also on the right hand side, albeit with a different sign if
$\left(  i,j\right)  $ is an inversion of $\sigma$. Thus, the products on both
sides agree up to a sign, which is precisely $\left(  -1\right)  ^{\ell\left(
\sigma\right)  }=\left(  -1\right)  ^{\sigma}$. \medskip

\textbf{(g)} This is \cite[Exercise 5.13 \textbf{(c)}]{detnotes}, and is a
particular case of Proposition \ref{prop.perm.sign.props} \textbf{(h)}.
\end{proof}

\begin{corollary}
\label{cor.perm.sign.hom}Let $n\in\mathbb{N}$. The map%
\begin{align*}
S_{n}  &  \rightarrow\left\{  1,-1\right\}  ,\\
\sigma &  \mapsto\left(  -1\right)  ^{\sigma}%
\end{align*}
is a group homomorphism from the symmetric group $S_{n}$ to the order-$2$
group $\left\{  1,-1\right\}  $. (Of course, $\left\{  1,-1\right\}  $ is a
group with respect to multiplication.)
\end{corollary}

This map is known as the \emph{sign homomorphism}.

\begin{proof}
[Proof of Corollary \ref{cor.perm.sign.hom} (sketched).]Proposition
\ref{prop.perm.sign.props} \textbf{(d)} shows that this map respects
multiplication (i.e., sends products to products). However, if a map between
two groups respects multiplication, then it is automatically a group
homomorphism. Thus, Corollary \ref{cor.perm.sign.hom} follows.
\end{proof}

\begin{definition}
\label{def.perm.even-odd}Let $n\in\mathbb{N}$. A permutation $\sigma\in S_{n}$
is said to be

\begin{itemize}
\item \emph{even} if $\left(  -1\right)  ^{\sigma}=1$ (that is, if
$\ell\left(  \sigma\right)  $ is even);

\item \emph{odd} if $\left(  -1\right)  ^{\sigma}=-1$ (that is, if
$\ell\left(  \sigma\right)  $ is odd).
\end{itemize}
\end{definition}

The sign and the \textquotedblleft parity\textquotedblright\ (i.e.,
evenness/oddness) of a permutation have applications throughout mathematics
(in the definition of determinants and the construction of exterior powers) as
well as in the solution of \emph{permutation puzzles} (such as Rubik's cube
and the $15$-game; see \cite[Chapters 7--8 and Theorem 20.2.1]{Mulhol21} for
example). Even permutations are also crucial in group theory, as they form a group:

\begin{corollary}
\label{cor.perm.altgp}Let $n\in\mathbb{N}$. The set of all even permutations
in $S_{n}$ is a normal subgroup of $S_{n}$.
\end{corollary}

This subgroup is known as the $n$\emph{-th alternating group} (commonly called
$A_{n}$). If $n\geq5$, then this group is a simple group (meaning that it has
no normal subgroups besides itself and the trivial group)\footnote{See, e.g.,
\url{https://groupprops.subwiki.org/wiki/Alternating_groups_are_simple} or
\cite[Theorem 10.3.4]{Goodman} for a proof.}, and this fact has been used by
Galois to prove that the general $5$-th degree polynomial equation cannot be
solved using radicals.

\begin{proof}
[Proof of Corollary \ref{cor.perm.altgp} (sketched).]The set of all even
permutations in $S_{n}$ is the kernel of the group homomorphism $S_{n}%
\rightarrow\left\{  1,-1\right\}  $ from Corollary \ref{cor.perm.sign.hom}.
Thus, it is a normal subgroup of $S_{n}$ (since any kernel is a normal subgroup).
\end{proof}

\begin{corollary}
\label{cor.perm.num-even}Let $n\geq2$. Then,%
\[
\left(  \text{\# of even permutations }\sigma\in S_{n}\right)  =\left(
\text{\# of odd permutations }\sigma\in S_{n}\right)  =n!/2.
\]

\end{corollary}

\begin{proof}
[Proof of Corollary \ref{cor.perm.num-even} (sketched).]The symmetric group
$S_{n}$ contains the simple transposition $s_{1}$ (since $n\geq2$). If
$\sigma\in S_{n}$, then
\begin{align*}
\left(  -1\right)  ^{\sigma s_{1}}  &  =\left(  -1\right)  ^{\sigma}%
\cdot\underbrace{\left(  -1\right)  ^{s_{1}}}_{\substack{=-1\\\text{(by
Proposition \ref{prop.perm.sign.props} \textbf{(b)},}\\\text{since }%
s_{1}=t_{1,2}\text{)}}}\ \ \ \ \ \ \ \ \ \ \left(  \text{by Proposition
\ref{prop.perm.sign.props} \textbf{(d)}}\right) \\
&  =-\left(  -1\right)  ^{\sigma}.
\end{align*}
Hence, a permutation $\sigma\in S_{n}$ is even if and only if the permutation
$\sigma s_{1}$ is odd. Hence, the map%
\begin{align*}
\left\{  \text{even permutations }\sigma\in S_{n}\right\}   &  \rightarrow
\left\{  \text{odd permutations }\sigma\in S_{n}\right\}  ,\\
\sigma &  \mapsto\sigma s_{1}%
\end{align*}
is well-defined. This map is furthermore a bijection (since $S_{n}$ is a
group). Thus, the bijection principle yields
\[
\left(  \text{\# of even permutations }\sigma\in S_{n}\right)  =\left(
\text{\# of odd permutations }\sigma\in S_{n}\right)  .
\]
Both sides of this equality must furthermore equal to $n!/2$, since they add
up to $\left\vert S_{n}\right\vert =n!$. This proves Corollary
\ref{cor.perm.num-even}. (See \cite[Exercise 5.4]{detnotes} for details.)
\end{proof}

As a consequence of Corollary \ref{cor.perm.num-even}, we see that%
\begin{equation}
\sum_{\sigma\in S_{n}}\left(  -1\right)  ^{\sigma}%
=0\ \ \ \ \ \ \ \ \ \ \text{for each }n\geq2.
\label{eq.cor.perm.num-even.sum-sign}%
\end{equation}
(Indeed, the sum $\sum_{\sigma\in S_{n}}\left(  -1\right)  ^{\sigma}$ can be
rewritten as
\[
\left(  \text{\# of even permutations }\sigma\in S_{n}\right)  -\left(
\text{\# of odd permutations }\sigma\in S_{n}\right)  ,
\]
since the addends corresponding to the even permutations $\sigma\in S_{n}$ are
equal to $1$ whereas the addends corresponding to the odd permutations
$\sigma\in S_{n}$ are equal to $-1$.)

We note that the sign can be defined not only for a permutation $\sigma\in
S_{n}$, but also for any permutation of any finite set $X$ (even if the set
$X$ has no chosen total order on it, as the set $\left[  n\right]  $ has).
Here is one way to do so:

\begin{proposition}
\label{prop.perm.sign.X}Let $X$ be a finite set. We want to define the sign of
any permutation of $X$.

Fix a bijection $\phi:X\rightarrow\left[  n\right]  $ for some $n\in
\mathbb{N}$. (Such a bijection always exists, since $X$ is finite.) For every
permutation $\sigma$ of $X$, set%
\[
\left(  -1\right)  _{\phi}^{\sigma}:=\left(  -1\right)  ^{\phi\circ\sigma
\circ\phi^{-1}}.
\]
Here, the right hand side is well-defined, since $\phi\circ\sigma\circ
\phi^{-1}$ is a permutation of $\left[  n\right]  $. Now: \medskip

\textbf{(a)} This number $\left(  -1\right)  _{\phi}^{\sigma}$ depends only on
the permutation $\sigma$, but not on the bijection $\phi$. (In other words, if
$\phi_{1}$ and $\phi_{2}$ are two bijections from $X$ to $\left[  n\right]  $,
then $\left(  -1\right)  _{\phi_{1}}^{\sigma}=\left(  -1\right)  _{\phi_{2}%
}^{\sigma}$.)

Thus, we shall denote $\left(  -1\right)  _{\phi}^{\sigma}$ by $\left(
-1\right)  ^{\sigma}$ from now on. We refer to this number $\left(  -1\right)
^{\sigma}$ as the \emph{sign} of the permutation $\sigma\in S_{X}$. (When
$X=\left[  n\right]  $, this notation does not clash with Definition
\ref{def.perm.sign}, since we can pick the bijection $\phi=\operatorname*{id}$
and obtain $\left(  -1\right)  _{\phi}^{\sigma}=\left(  -1\right)
^{\operatorname*{id}\circ\sigma\circ\operatorname*{id}\nolimits^{-1}}=\left(
-1\right)  ^{\sigma}$.) \medskip

\textbf{(b)} The identity permutation $\operatorname*{id}:X\rightarrow X$
satisfies $\left(  -1\right)  ^{\operatorname*{id}}=1$. \medskip

\textbf{(c)} We have $\left(  -1\right)  ^{\sigma\tau}=\left(  -1\right)
^{\sigma}\cdot\left(  -1\right)  ^{\tau}$ for any two permutations $\sigma$
and $\tau$ of $X$.
\end{proposition}

\begin{proof}
[Proof of Proposition \ref{prop.perm.sign.X} (sketched).]This all follows
quite easily from Proposition \ref{prop.perm.sign.props} \textbf{(d)}. See
\cite[Exercise 5.12]{detnotes} for a detailed proof.
\end{proof}
